\documentclass[runningheads]{llncs}

\usepackage[T1]{fontenc}
\usepackage{amsmath,amssymb}
\usepackage{booktabs}
\usepackage{listings}
\usepackage{xcolor}
\usepackage[hidelinks]{hyperref}

\setlength{\emergencystretch}{1em}

\newcommand{\fpdiv}{\mathop{\mathsf{fp.div}}}
\newcommand{\isnan}{\mathsf{NaN}}
\newcommand{\isinf}{\mathsf{Inf}}
\newcommand{\iszero}{\mathsf{Zero}}
\newcommand{\signbit}{\mathsf{Neg}}
\newcommand{\proxy}{\alpha_d}
\newcommand{\true}{\mathsf{true}}
\newcommand{\false}{\mathsf{false}}

\lstdefinestyle{plan}{
  basicstyle=\ttfamily\small,
  columns=fullflexible,
  keepspaces=true,
  frame=single,
  rulecolor=\color{black!25},
  xleftmargin=0.5em,
  xrightmargin=0.5em,
  aboveskip=0.7em,
  belowskip=0.7em
}

\title{Term-Level Abstraction and Classification Refinement for
Floating-Point Division in Bitwuzla}

\titlerunning{FP Division Abstraction in Bitwuzla}

\author{Anonymous Author(s)}
\authorrunning{Anonymous Author(s)}
\institute{Anonymous Institution\\
\email{anonymous@example.org}}

\begin{document}

\maketitle

\begin{abstract}
This design note proposes the first floating-point extension of Bitwuzla's
term-level abstraction module.  The initial scope is deliberately restricted
to floating-point division.  Each selected \(\fpdiv\) term is replaced by a
fresh, valid floating-point proxy.  Candidate models are checked against the
concrete division semantics, and refinement starts with inexpensive
classification lemmas for NaN, infinity, zero, and sign.  If classification
refinement cannot eliminate an inconsistent candidate, the first prototype
falls back to an exact semantic link and the existing SymFPU encoding.  This
preserves soundness and completeness while isolating the main research
question: how much of the full divider can be avoided by class-first
refinement?

\keywords{SMT solving \and floating-point arithmetic \and abstraction
\and refinement \and bit-blasting}
\end{abstract}

% CAV uses the Springer LNCS format, whose top-level units are sections.
% Sections 1 and 2 correspond to the two requested design chapters.

\section{Top-Level Design}
\label{sec:top-level}

\subsection{Scope and design goals}

The first milestone adds term-level abstraction for \emph{only}
\(\mathsf{FP\_DIV}\).  Let
\[
  d = \fpdiv(r,x,y),
\]
where \(r\) is a rounding-mode term and \(x,y\) have the same floating-point
format.  After rewriting, the abstraction module replaces every occurrence of
\(d\) by a shared proxy \(\proxy\) of the same floating-point type:
\[
  \varphi[d] \quad\longmapsto\quad \widehat{\varphi}[\proxy].
\]
The proxy initially has no division semantics.  It is constrained only to be
a valid floating-point value.  Consequently,
\(\widehat{\varphi}\) over-approximates the original formula.

The milestone has four goals.
\begin{enumerate}
  \item Reuse the existing \texttt{AM\_ABSTRACT} representation and the
        abstraction/refinement loop introduced for bit-vectors
        \cite{niemetz2024abstraction}.
  \item Prevent the original \(\fpdiv\) term from being word-blasted until an
        exact fallback is required.
  \item Add operation-specific classification lemmas before any full SymFPU
        division encoding.
  \item Preserve the result of Bitwuzla with abstraction disabled: no
        approximate SAT answer is permitted.
\end{enumerate}

The following features are explicitly outside the first milestone:
\(\mathsf{FP\_ADD}\), \(\mathsf{FP\_MUL}\), \(\mathsf{FP\_FMA}\),
\(\mathsf{FP\_SQRT}\), predicates with Boolean result, conversions with
bit-vector result, and partial operators such as \(\mathsf{FP\_MIN}\) and
\(\mathsf{FP\_MAX}\).  Ground divisions are expected to disappear through
constant folding before abstraction.

\subsection{Semantic contract}

The implementation must maintain the following invariants.

\paragraph{Typed proxy.}
For each rewritten node \(d\), create exactly one proxy \(\proxy\), preserve
DAG sharing, and give \(\proxy\) the type of \(d\).  A floating-point proxy is
treated as a symbolic leaf by the FP word-blaster.  Its unpacked
NaN/Inf/Zero/sign/exponent/significand fields must satisfy the same validity
constraint as an ordinary symbolic FP constant.

\paragraph{No accidental expansion.}
The word-blaster must stop at \(\proxy\); it must not descend through the
internal child of \texttt{AM\_ABSTRACT} to \(d\).  Classification lemmas may
mention \(r,x,y,\proxy\), but must not mention \(d\).  Otherwise merely
classifying the result could construct the complete symbolic divider.

\paragraph{Exact candidate validation.}
For every SAT candidate, obtain the model values
\(v_r,v_x,v_y,v_{\proxy}\) and evaluate
\[
  v_d = \fpdiv(v_r,v_x,v_y)
\]
with Bitwuzla's existing FP evaluator.  The abstraction is consistent only if
\(v_{\proxy}=v_d\) under SMT equality.  In particular, this equality
distinguishes \(+0\) from \(-0\) and uses the SMT treatment of NaN; it is
\emph{not} the IEEE predicate \(\mathsf{fp.eq}\).

\paragraph{Complete fallback.}
Classification lemmas alone do not determine a finite quotient.  If the
candidate remains inconsistent after all applicable class refinements, the
first implementation adds
\[
  \proxy = \fpdiv(r,x,y).
  \tag{Full-Link}
\]
The lemma must be passed through the ``do not abstract this lemma'' path so
that the right-hand-side division is word-blasted rather than abstracted
again.  This fallback bounds the refinement loop and restores the current
SymFPU semantics \cite{symfpu}.

\subsection{Solver pipeline}

The proposed data flow is as follows.

\begin{enumerate}
  \item \textbf{Rewrite and select.}  Process the formula bottom-up.  If a
        rewritten node has kind \(\mathsf{FP\_DIV}\), contains symbolic data,
        and the new option is enabled, replace it by \(\proxy\).
  \item \textbf{Build the abstract FP formula.}  The existing FP solver
        word-blasts surrounding predicates and equalities, but treats
        \(\proxy\) as a valid symbolic FP leaf.  The arithmetic divider for
        \(d\) is absent from this encoding.
  \item \textbf{Obtain a candidate.}  The BV/SAT engine returns an assignment;
        the current solve order then lets the FP solver establish valid proxy
        fields before the abstraction module checks \(\proxy\).
  \item \textbf{Check exact semantics.}  Evaluate \(d\) on the candidate
        operand values.  If \(v_{\proxy}=v_d\), this abstraction is consistent.
  \item \textbf{Refine class first.}  On a mismatch, test the classification
        lemmas from Sect.~\ref{sec:class-lemmas} and add a violated lemma.
  \item \textbf{Fall back if necessary.}  If no class lemma rejects the
        inconsistent candidate, add Full-Link.  Later work may insert value,
        exponent, or significand refinements before this step.
\end{enumerate}

The first implementation should retain the existing ordering
\[
  \text{BV/SAT solve} \to \text{FP check} \to
  \text{abstraction check} \to \text{array/UF checks}.
\]
Reordering is unnecessary if FP proxies are valid word-blaster leaves and
their model values can be recovered through \texttt{FpSolver::value()}.

\subsection{Required code changes}

Table~\ref{tab:changes} summarizes the expected change surface.  Keeping the
FP-specific checker separate from the current binary-BV checker avoids
hard-coded assumptions about arity and bit-vector width.

\begin{table}[t]
\centering
\caption{Planned implementation changes.}
\label{tab:changes}
\small
\begin{tabular}{p{0.25\textwidth}p{0.66\textwidth}}
\toprule
Component & Planned change \\
\midrule
Options
  & Add an experimental, default-off \texttt{abstraction-fp-div} option under
    the existing abstraction master switch.  Do not reuse
    \texttt{abstraction-bv-size}. \\
Abstraction selection
  & Extend \texttt{AbstractionModule::abstract()} with an FP-specific branch
    for \texttt{Kind::FP\_DIV}.  Reuse the process cache to preserve sharing. \\
FP refinement checker
  & Add \texttt{check\_term\_abstraction\_fp\_div()}; read one RM and two FP
    operands, perform exact candidate evaluation, then use the class-first
    schedule. \\
FP word-blaster
  & Recognize FP-typed \texttt{AM\_ABSTRACT} as a symbolic FP leaf, add its
    validity constraint, and do not traverse its hidden original child. \\
Model values
  & Permit \texttt{FpSolver::value()} to recover the packed value of an
    FP-typed abstraction proxy. \\
Lemma handling
  & Ensure class lemmas contain no original \texttt{FP\_DIV}.  Cache and
    backtrack lemmas like existing abstraction lemmas.  Mark Full-Link as
    non-abstractable. \\
Statistics
  & Count abstracted FP divisions, candidate checks, each class-lemma kind,
    exact matches, and full-link fallbacks. \\
Tests
  & Add unit tests for proxy typing/word-blasting and differential regressions
    comparing abstraction on and off. \\
\bottomrule
\end{tabular}
\end{table}

\subsection{Refinement algorithm}

The control structure should be explicit about the distinction between a
class mismatch and a numeric mismatch.

\begin{lstlisting}[style=plan,caption={Class-first checking for one abstracted FP division.}]
check_fp_div(proxy a, term div(r, x, y)):
    vr, vx, vy, va := model_value(r, x, y, a)
    expected       := eval_fp_div(vr, vx, vy)
    if smt_equal(va, expected):
        return CONSISTENT

    for lemma in [NAN_POS, NAN_NEG, INF_SPECIAL,
                  ZERO_SPECIAL, SIGN]:
        if lemma is false in the candidate:
            add(lemma(r, x, y, a))
            return REFINED

    add_no_abstract(a = fp.div(r, x, y))
    return REFINED_BY_FULL_LINK
\end{lstlisting}

The normal mode adds the first violated class lemma, matching the current lazy
abstraction strategy.  An eager mode may add all violated class lemmas in one
iteration.  Exact candidate evaluation remains mandatory in either mode.

\subsection{Correctness argument}

The abstract formula is an over-approximation because any model of the
original formula extends to a model of the abstraction by assigning
\(\proxy=d\).  Every class lemma in Sect.~\ref{sec:class-lemmas} is a logical
consequence of IEEE/SMT floating-point division \cite{IEEE754}; therefore, an
UNSAT result after adding these lemmas implies UNSAT of the original formula.

A SAT result is returned only after every active proxy passes exact candidate
validation.  Replacing each proxy value by the corresponding evaluated
division value then preserves all surrounding constraints, yielding a model
of the original formula.  Finally, Full-Link reduces unresolved cases to the
existing complete word-blasted encoding.  Thus the initial implementation is
sound and complete even if the class lemmas provide no performance gain.

\subsection{Milestones and evaluation}

The recommended implementation order is:
\begin{enumerate}
  \item FP proxy creation, validity, and model-value recovery;
  \item abstraction with immediate Full-Link, establishing an end-to-end
        correctness baseline;
  \item the NaN lemmas from Sect.~\ref{subsec:nan};
  \item special infinity, zero, and sign lemmas;
  \item statistics and AIG/CNF measurements; and
  \item only after the measurements, finer exponent or significand lemmas.
\end{enumerate}

For each benchmark, measure wall time, number of SAT calls, number and kind of
refinements, number of full-link fallbacks, and reachable AIG/CNF size.  The
most informative encoding comparison is: full \(\fpdiv\), a class predicate
obtained through the full SymFPU result, and the direct class lemma proposed
below.

\section{Design of Classification Lemmas}
\label{sec:class-lemmas}

\subsection{Notation and source semantics}

For a valid floating-point term \(t\), define
\[
\begin{array}{lll}
  \isnan(t)  &\text{iff }t\text{ is NaN},&
  \isinf(t)  \text{ iff }t\text{ is infinite},\\
  \iszero(t) &\text{iff }t\text{ is }+0\text{ or }-0,&
  \signbit(t)\text{ iff }t\text{ is negative}.
\end{array}
\]
The implementation should use Bitwuzla's existing floating-point
classification nodes for NaN, infinity, zero, and negativity.  SymFPU defines
NaN as neither positive nor negative, so sign lemmas are guarded against NaN.

The pinned SymFPU implementation first constructs and rounds an arithmetic
quotient, then applies three prioritized bypasses:
\[
  \mathsf{if}\ C_N\ \mathsf{then}\ \mathrm{NaN}\
  \mathsf{else\ if}\ B_I\ \mathsf{then}\ \mathrm{Inf}\
  \mathsf{else\ if}\ B_Z\ \mathsf{then}\ \mathrm{Zero}\
  \mathsf{else}\ \mathrm{roundedQuotient}.
\]
This structure appears in
\path{subprojects/symfpu/core/divide.h}.  It is the semantic source for the
class lemmas below.

Let \(x\) be the dividend and \(y\) the divisor.  Introduce the operand-only
conditions
\begin{align}
  C_N ={}& \isnan(x) \lor \isnan(y)
       \lor (\isinf(x) \land \isinf(y))
       \lor (\iszero(x) \land \iszero(y)),
       \label{eq:cn}\\
  B_I ={}& (\neg\iszero(x) \land \iszero(y))
       \lor (\isinf(x) \land \neg\isinf(y)),
       \label{eq:bi}\\
  B_Z ={}& (\neg\isinf(x) \land \isinf(y))
       \lor (\iszero(x) \land \neg\iszero(y)).
       \label{eq:bz}
\end{align}
To mirror the prioritized ITE exactly, define
\begin{align}
  C_I & = \neg C_N \land B_I,
       \label{eq:ci}\\
  C_Z & = \neg C_N \land \neg B_I \land B_Z.
       \label{eq:cz}
\end{align}
The \(\neg B_I\) conjunct is redundant for valid non-NaN operand classes, but
retaining it documents and preserves the SymFPU branch priority.

The rounding mode does not occur in Eqs.~\eqref{eq:cn}--\eqref{eq:cz}:
these conditions describe only the special-value bypasses.  Rounding is
relevant to finite overflow and underflow, which is why some lemmas below are
one-directional.

\subsection{Exact NaN refinement}
\label{subsec:nan}

Division produces NaN exactly when an operand is NaN, both operands are
infinite, or both operands are zero.  Rounding a non-special quotient does not
create NaN.  Hence the strongest cheap class fact is
\[
  \isnan(\proxy) \leftrightarrow C_N.
  \tag{NAN}
\]
For lazy candidate-driven refinement, split the equivalence into two lemmas:
\begin{align}
  L_{N+}:&\quad C_N \Longrightarrow \isnan(\proxy),
       \label{lem:nan-pos}\\
  L_{N-}:&\quad \isnan(\proxy) \Longrightarrow C_N.
       \label{lem:nan-neg}
\end{align}
This permits the solver to add only the violated direction.  For example,
\(0/0\), \(\infty/\infty\), and any operation with a NaN operand violate
\(L_{N+}\) whenever the proxy is assigned a non-NaN value.  Conversely,
\(L_{N-}\) eliminates a spurious NaN proxy for ordinary finite operands.

\subsection{Special infinity refinement}

The infinity bypass covers a nonzero value divided by zero and infinity
divided by a non-infinite value.  The corresponding sound lemma is
\[
  L_I:\quad C_I \Longrightarrow \isinf(\proxy).
  \tag{INF-SPECIAL}
\]
This lemma must \emph{not} be strengthened to an equivalence.  A finite,
nonzero quotient may overflow during rounding and become infinity even when
\(C_I\) is false.  The exact overflow condition depends on exponent,
significand, sign, and rounding mode and belongs to a later refinement stage.

\subsection{Special zero refinement}

The zero bypass covers a non-infinite value divided by infinity and zero
divided by a nonzero value.  The corresponding lemma is
\[
  L_Z:\quad C_Z \Longrightarrow \iszero(\proxy).
  \tag{ZERO-SPECIAL}
\]
Again, the converse is unsound: a finite nonzero quotient may underflow and
round to signed zero while \(C_Z\) is false.  The SymFPU rounder explicitly
selects zero or the minimum subnormal according to the rounding mode and sign.

\subsection{Sign refinement}

For every non-NaN division result, IEEE division uses the exclusive-or of the
operand signs.  This remains true for normal, subnormal, infinite, and signed
zero results.  Since \(\signbit(t)\) is false for NaN in the SMT-LIB
classification semantics, use the guarded lemma
\[
  L_S:\quad
  \neg\isnan(\proxy) \Longrightarrow
  \bigl(\signbit(\proxy)
  \leftrightarrow
  (\signbit(x)\oplus\signbit(y))\bigr).
  \tag{SIGN}
\]
When the guard holds, Eq.~\eqref{eq:cn} and the NaN lemmas ensure that neither
operand is NaN, so \(\signbit(x)\) and \(\signbit(y)\) expose their actual sign
bits.  The rounding mode is irrelevant to this lemma.

\subsection{Candidate-driven lemma selection}

Table~\ref{tab:selection} defines the initial lazy policy.  Lemmas are tested
in table order.  Earlier lemmas establish cleaner class assignments for later
checks.

\begin{table}[t]
\centering
\caption{Class-first refinement decisions for a candidate model.}
\label{tab:selection}
\small
\begin{tabular}{p{0.43\textwidth}p{0.21\textwidth}p{0.25\textwidth}}
\toprule
Observed candidate condition & Add & Purpose \\
\midrule
\(C_N\) and \(\neg\isnan(\proxy)\)
  & \(L_{N+}\) & Force a required NaN. \\
\(\neg C_N\) and \(\isnan(\proxy)\)
  & \(L_{N-}\) & Reject a spurious NaN. \\
\(C_I\) and \(\neg\isinf(\proxy)\)
  & \(L_I\) & Enforce a special infinity. \\
\(C_Z\) and \(\neg\iszero(\proxy)\)
  & \(L_Z\) & Enforce a special signed zero class. \\
Non-NaN proxy and incorrect sign XOR
  & \(L_S\) & Enforce the division sign. \\
No class violation, but \(v_{\proxy}\ne v_d\)
  & Full-Link & Restore complete division semantics. \\
\bottomrule
\end{tabular}
\end{table}

All conditions are evaluated using the current model before building a lemma.
The instantiated formulas themselves remain symbolic and globally valid.  A
lemma cache prevents duplicates across repeated candidates.

\subsection{Cost discipline}

The central implementation rule is that a class lemma may not construct the
original division term.  In particular, it must not contain
\[
  \fpdiv(r,x,y).
\]
For example, the right-hand side of \(L_{N+}\) must be constructed directly
from classifications of \(x\) and \(y\), rather than as
\(\isnan(\fpdiv(r,x,y))\).  The latter calls SymFPU's complete
\texttt{divide()} constructor before selecting the NaN field.  Even if dead
quotient components are later excluded from the reachable SAT cone, their DAG
construction cost is paid and the intended abstraction boundary becomes
unclear.

Direct lemmas use only classifications already supported by the FP
word-blaster.  Their reachable cones contain operand/proxy validity and a
small Boolean network; they do not contain fixed-point significand division,
normalization, or rounding.

\subsection{Validation matrix}

Table~\ref{tab:matrix} is the minimum semantic test matrix.  Each row must be
tested with all sign combinations where meaningful and with symbolic as well
as concrete rounding modes.

\begin{table}[t]
\centering
\caption{Minimum tests for classification refinement.}
\label{tab:matrix}
\small
\begin{tabular}{lll}
\toprule
Operands & Required result class & Primary lemma \\
\midrule
NaN / any, any / NaN & NaN & \(L_{N+}\) \\
\(+\infty/-\infty\) & NaN & \(L_{N+}\) \\
\(+0/-0\) & NaN & \(L_{N+}\) \\
finite nonzero / zero & signed infinity & \(L_I,L_S\) \\
infinity / finite nonzero & signed infinity & \(L_I,L_S\) \\
finite / infinity & signed zero & \(L_Z,L_S\) \\
zero / finite nonzero & signed zero & \(L_Z,L_S\) \\
large finite / tiny finite & possible rounded infinity & no converse \(L_I\) \\
tiny finite / large finite & possible rounded zero & no converse \(L_Z\) \\
ordinary finite / finite & exact finite quotient & Full-Link initially \\
\bottomrule
\end{tabular}
\end{table}

Differential regression tests should solve every case with FP division
abstraction both enabled and disabled, compare SAT/UNSAT, and validate models.
Formats should include Float16, Float32, Float64, and Float128.  If experimental
formats are enabled, exhaustive tiny formats are particularly useful for
checking the class lemmas and the exact candidate evaluator.

\subsection{Acceptance criteria for the first implementation}

The FP-division milestone is complete only when all of the following hold:
\begin{enumerate}
  \item no \(\mathsf{FP\_DIV}\) arithmetic circuit is reachable before a
        refinement requires it;
  \item all class lemmas are validated against exhaustive small-format values
        or an independent abstraction-off solver;
  \item an accepted SAT candidate has exact proxy values, not merely matching
        classes;
  \item the Full-Link fallback cannot be re-abstracted;
  \item model production, incremental push/pop, unsat cores, and nested/shared
        division terms remain correct; and
  \item abstraction-off behavior and performance are unchanged.
\end{enumerate}

\bibliographystyle{splncs04}
\bibliography{references}

\end{document}
